% Generated by D:/tools/scratch_qdd/mmrec/render.py from data_b.py - edit the data, not this file.
\subsection{B. Astra 대체 --- 학습할 수 있는 크기의 상위 모델로}\label{rec:B}
\noindent\textbf{한 줄:} 목표는 학습 가능한 크기의 오픈 VLM이 Astra 자리를 맡는 것(§96). 영샷은 gpt-5.2 대리 0/4·8B 0/5로 인식에서 막혔고, 무료 시뮬 참값으로 LoRA 학습한 8B는 4/4(접근 오차 125.5 $\rightarrow$ 2.6 mm). 최종 상위는 35B, 8B는 임시. 손으로 준 정보 없이도 되는지(E-STRIP8)와 거리를 모델이 아는 두 방법(E-PT)은 등록 뒤 진행 중.

\paragraph{B1. 배경: Jev를 못 쓰게 됨 $\rightarrow$ 로컬 VLM 선택기}
{\small 09-24 18:44 -- 19:10 KST}
\begin{itemize}
  \item \textbf{사용자} (ul 46, 09-24 18:44 KST): \intent{jev는 지금 사용이 불가능한상황임을 확인했어 애초에 텍스트 기반의 api 활성화 기반자체가 우리랑 엮기에는 약간은 부족할 수 있겠단 생각을 해서 jev를 대체할 무언가 방법을 찾아서 해볼 필요가 있을 것 같아}
  \item \textbf{사용자} (ul 47, 19:10 KST): \intent{아스트라써야해 알지?}
  \item \textbf{그래서(결정):} 정본 §44(09-24 18:55 KST): 하위 결정기 = 로컬 VLM typed 선택기 Jev-L(Qwen3-VL, vLLM, 보기 토큰 로그 확률). 사전 시험: 4B 텍스트 + 머리캠 동시 4개 p95 0.261 s. Astra는 위층 계획기로 유지(ul 47 선택).
  \item \textbf{참고:} 이 하위 자리는 이후 융합 VLA(Qwen3-VL-4B + 행동 전문가, §51·§58)로 흡수되었고, 이 줄기(B)는 위층 Astra 자리를 무엇이 대신할지를 다룬다.
\end{itemize}

\paragraph{B2. 선생-학생 $\rightarrow$ 목표 재정의: 학습 가능한 크기로 Astra 대체}
{\small 09-26 19:4x -- 20:00 KST}
\begin{itemize}
  \item \textbf{사용자} (ul 108, 19:4x KST): \intent{선생으로 활약을 잘한다면 꽤 괜찮을거라고 보거든? Llm의 단점을 해결해야해 느린 답변 Vla의 낮은 정확성}
  \item \textbf{사용자} (ul 108, 19:4x KST): \intent{둘중 옳은 방식으로 해야지}
  \item \textbf{그래서(결정):} 정본 §96(047ebc8, 19:57): 시뮬 선생 데이터(E-TEACH-0: 학생 VLA가 시뮬 폐루프를 몰고 선생이 라벨)는 본학습 자리만, 실물 동기 선생 조종(E-TEACH-R)은 추가학습 자리(실물 접근이 없어 계획만). 선행 조건 = 런타임 멈춤 고리 수정(C7).
  \item \textbf{사용자} (ul 109, 20:0x KST): \intent{나는 학습할 수 있는 llm크기에서 아스트라를 대체하는게 목적이야}
  \item \textbf{그래서(결정):} §96 개정 1(ce9c567, 20:00): 주 학생 = 학습 가능 크기 오픈 VLM(Qwen3.5-27B·35B-A3B 급, H200 1--2장 LoRA; 확장 122B-A10B QLoRA), Astra는 증류 선생(E-TEACH-L). 성공 = DEV 짝 40에서 Astra 단독 비열등(여유 3/40), 지연 p50 $\leq$ Astra의 1/3, API 0원. 양은 무료 참값 모델(\texttt{astra\_\allowbreak{}solo/\allowbreak{}truth.\allowbreak{}py}), Astra는 핵심 상태만(단계 $\leq$ 8,000원); 팔 L-T(참값)·L-A(Astra)·L-TA. VLA DAgger는 보조 트랙 E-TEACH-V.
\end{itemize}

\paragraph{B3. 영샷 오픈급 모델로 바로 대체? --- gpt-5.2 대리 0/4}
{\small 09-26 19:45 -- 20:17 KST}
\begin{itemize}
  \item \textbf{해 봄:} 조사 4b91f98(19:45): 목표 오픈 VLM Qwen3.5-397B-A17B(Apache-2.0, 4 H200 FP8로 호스팅 가능), OpenAI 대리 gpt-5.2(같은 Qwen 카드 표 기준). 내려받던 부분 파일 96 GB는 사용자 정정으로 멈추고 삭제. 사전 등록 5c021c6(19:49, 사용자 `허용': gpt-5.2 $\times$ 4편, 상한 3,000원): \texttt{astra-\allowbreak{}solo@v2} 그대로 모델 id만 바꿈, effort low, 조기 종료 20호출/60 s.
  \item \textbf{참고:} 등록 변경 1(f37051c, 20:06): 사용자 정정 --- 목표는 우리가 미세조정할 수 있는 크기(Qwen3.5-27B), 대리 gpt-5-mini 추가. gpt-5-mini는 7호출(23.3원) 뒤 사용자 지시로 멈춤(부분 편, 판정에 안 씀).
  \item \textbf{결과:} 4fbf4e7(20:17): gpt-5.2 0/4, 파지 + 들기 0/4, 무효 0/79. 접근 목표 xy 편 중앙 104--207 mm(같은 판 Astra 2.6--48.9 mm) --- 격자 깊이 방향 x를 12--19 cm 멀게 읽음. 16.4원/호출, 지연 p50 7.2 s. 1,299.1원(+ mini 23.3원).
  \item \textbf{그래서(결정):} 판정 PERCEPTION: 영샷으로는 격자 위치 읽기를 못 한다 $\rightarrow$ 학습 가능 크기 오픈 VLM도 영샷 대체는 기대하기 어렵다; E-TEACH-L의 출발점은 `인식 0', 교사 데이터는 접근 목표 좌표(인식)를 가장 많이 가르쳐야 한다. 같은 인터페이스의 무료 Qwen3-VL-8B 영샷도 0/5(A8 사전 실행).
\end{itemize}

\paragraph{B4. 로봇 자기 정보 말고는 전부 자동화 + 모델 수 최소화 + 깊이 가정}
{\small 09-26 20:3x -- 20:47 KST}
\begin{itemize}
  \item \textbf{사용자} (ul 110, 20:3x KST): \intent{음 이 뭐랄까 로봇의 형상만 보고도 로봇의 엔드포인트를 알아서 알아내고 그것이 행하는 방향을 알아내는걸 자동화 하는걸 할필요가 있어보임 지금 우리가 준 정보라는 것 자체들 전부다 그리퍼 형상, 카메라,범위 같은 로봇이 자기가 스스로 할 수 있는것 외에는 다 자동화를 할 수 있어야해 근데 이게 뭐 여러모델을 이어붙이는건 좋은데 그 모델을 최소화 해야하는 게 좋은 목표인거고}
  \item \textbf{사용자} (ul 110, 20:3x KST): \intent{댑스를 쓴다는걸 가정을 해야겠다}
  \item \textbf{그래서(결정):} 정본 §97(36a7309, 20:47): \texttt{astra-\allowbreak{}solo@v2} 인벤토리 R 16 · P 5 · L 3 · X 3 · K 1. 로봇 자기 정보(R)는 URDF·관절 상태·camera\_info·자기 시험에서 코드가 자동 생성해 계속 준다; 장면(P)은 깊이 + VLM 점 + 코드(RANSAC 평면·평면 위 군집·역투영); 요령(L)은 학습; X는 절제로 판단. 시각 자기 인식은 상위 VLM의 점 출력(시뮬 FK 무료 라벨) + FK 대조. 깊이: 시뮬 렌더 깊이, 실물 머리 ZED Mini 스테레오 계산·손목 D405. 모델 목표 2개(상위 VLM + VLA).
  \item \textbf{다음:} 절제 사다리 A0(무료) $\rightarrow$ A1(무료) $\rightarrow$ A2(유료 $\leq$ 8,000원) --- 모두 [예정], 결과 없음.
\end{itemize}

\paragraph{B5. 가설 H-L --- 상관관계만 학습한 상위 LLM이 더 빠르고 비슷하거나 낫다}
{\small 09-26 21:44 -- 21:50 KST}
\begin{itemize}
  \item \textbf{사용자} (ul 112, 21:44 KST): \intent{사실 실시간으로 아스트라를 키기에는 문제가 좀 있다는 생각이 있어 아스트라가 주역인데 아스트라가 실시간적으로 참여하기에는 너무 길다 감독을 하기에는 시간이 너무 길다는 문제가 있음 따라서 상위 llm이 아스트라보다는 성능이 낮지만 이 엔드포인트와 로봇의 상관관계 이동의 상관관계만을 학습한 것이 더 빠르고 성ㅇ능은 비슷 오히려 더 좋을 것이다가 가정이야}
  \item \textbf{사용자} (ul 113, 21:50 KST): \intent{이런 말두ㅡㄹ도 아직 확정은 아니지만 논문에 미리 업뎃해두고 기억해줘야헤}
  \item \textbf{그래서(결정):} 정본 §96 보충 1(e4156f8, 21:48): H-L 등록, 측정 기준 (a) 상위 호출 지연 p50 $\leq$ 2 s·Astra의 1/3 이하, (b) DEV 짝 폐루프 성공이 Astra 단독과 3/40 이내, (c) 접근 xy 오차 $\leq$ Astra(중앙 약 14 mm)·7--10 cm 오표적 감지 > Astra 0/13. 전제 표 P13. 논문 수시 갱신(054f250, 21:48, 주황 표시).
  \item \textbf{참고:} 당시 근거: gpt-5.2 대리 0/4(104--207 mm), Qwen3-VL-8B 0/5, Astra 자신도 7--10 cm 드리프트 0/13; Qwen3.5-35B-A3B 지연 약 1--2 s는 추정(실측 아님).
\end{itemize}

\paragraph{B6. 35B 전에 8B로 경향부터 --- E-TEACH-L8, 4/4}
{\small 09-26 21:0x -- 23:53 KST}
\begin{itemize}
  \item \textbf{사용자} (ul 117, 21:0x KST): \intent{학습을하면 35b에서 될것이다하는 가설이 진짜 너무 좋긴한데 이게 학습시간이 너무 오래걸리기 때문에 바로 35로가기는 너무 문제가 큼 그래서 8B를 학습을 해서 경향성이 보이는지만 파악을 하면 더 좋을 것 같기는함}
  \item \textbf{해 봄:} 사전 등록 72efee3(21:20), 변경 1 e7960de(21:49, 학습 전: 미세 배치 8 $\times$ 누적 2). Qwen3-VL-8B LoRA SFT, 시뮬 참값 라벨을 \texttt{astra-\allowbreak{}solo@v2} 형식 그대로(DART식 교란 포함). TRAIN 300편 $\rightarrow$ 조종 2,496행(반복 가중 뒤 4,025) + 인식 보조 QA 2,496 = 에폭당 6,521; LoRA r16 2 에폭.
  \item \textbf{결과:} c253ba1(23:48): 오프라인 DEV 305 스냅숏 접근 목표 xy 중앙 영샷 125.5 $\rightarrow$ 미세조정 2.6 mm(대상 위 단계 158.8 $\rightarrow$ 7.2, 첫 호출 163.4 $\rightarrow$ 8.1), 그리퍼 행동 정확도 0.32 $\rightarrow$ 0.95, 유효 JSON 1.00. 폐루프 DEV 4판: 미세조정 4/4(편당 6호출, 성공까지 11.5--11.9 s, 접근 편 중앙 7.7--19.0 mm) 대 영샷 0/4. 학습 1.39 GPU-h(약 4.4k 토큰/s), 지연 p50 3.5 s, 유료 0원.
  \item \textbf{그래서(결정):} 판정 TREND·STRONG(오프라인) $\rightarrow$ WORKS(폐루프). DAgger는 안 함(GPU 조건 초과·이미 천장). 뜻: 영샷 실패는 크기보다 과제 데이터가 없던 문제 쪽; 8B가 이 과제에서 천장이라 35B의 크기 효과는 가르지 못한다. 논문 수시 갱신(35c9fbb, 23:53).
  \item \textbf{다음:} 더 어려운 평가(OOD 높이·물체·방해물, 새 과제, 실물 프레임)와 다양한 데이터(E5·E6), 35B-A3B 처리량 스모크. 한계: n = 4, 참값 라벨(Astra 라벨 아님), 과제 하나, 프롬프트에 손으로 쓴 정보가 있음($\rightarrow$ B9).
\end{itemize}

\paragraph{B7. m 좌표 없이 조종 $\rightarrow$ 점 찍고 행동 + 거리를 모델이 아는 두 방법(E-PT)}
{\small 09-26 23:4x -- 09-27 00:50 KST}
\begin{itemize}
  \item \textbf{사용자} (ul 118, 23:4x--23:5x KST): \intent{로봇 좌표로 x, y, z가 몇 m인지 없이 할방법을 관련 서적 다뒤져서 찾아봐보 ㅏ뎁스여도되고 이미지여도 되고 다괜찮아}
  \item \textbf{사용자} (ul 118, 23:4x--23:5x KST): \intent{로봇 좌표 x, y, z를 m 단위로 말하지 않고 조종하는 방법도 좋지만 어떻게 그걸 추정할 수 있는지에 대한 방법론도 괜ㅊ낳아}
  \item \textbf{결과:} 조사 1d81290(00:01): 상위 VLM은 m 숫자 대신 영상 위에서 고른다(점 $\rightarrow$ 물체 ID $\rightarrow$ 상대 방향), m 값은 코드가 깊이·카메라 자기 정보로 계산. 근거: SPF 같은 백본에서 점 100 \% 대 글 수치 7 \% 대 후보 선택 40 \%; 직접 metric은 SpatialRGPT $\pm$25 \% 성공 41 \%, CA-MLLM 카메라 크기 $\times$ 0.8에 F1 46.5 $\rightarrow$ 25.8. 추천 1 점 + 역투영 + 평면 위 군집, 2 물체 ID 선택, 3 상대 방향 $\times$ 크기 등급 + 손목 서보.
  \item \textbf{사용자} (ul 120, 23:4x KST): \intent{"그리퍼가 어디로 가야 하나"를 이미지 위에서 배우는 방식으로 해봐야지 지금 우리가 못하는게 뭔데 부족한게 뭐야?}
  \item \textbf{사실 확인:} 부족했던 것: (i) \texttt{astra\_\allowbreak{}solo}에 픽셀 목표 명령 없음, (ii) \texttt{astra\_\allowbreak{}solo} 월드가 깊이를 끔, (iii) 우리 데이터에 픽셀 라벨 없음(자체 시뮬은 투영으로 무료), (iv) 공개 MolmoAct식 변환은 다른 에이전트 몫.
  \item \textbf{사용자} (ul 121, 09-27 00:0x KST): \intent{근데 거리정보나 이런거 없이도 되면 너무 더 좋은 결과가 될 수 있을것 같은데 두가지 방법론이야 없게 해서 기존데이터를 그냥 사용하돼 그 거리 정보를 잘 알 수 있게 하던지 아니면 거리를 크리를 알 수 있게 하던지}
  \item \textbf{사용자} (ul 122, 00:1x KST): \intent{2 방법 모두 아는 방법으로 가야겠어 하자 해보자}
  \item \textbf{그래서(결정):} 정본 §97 보충 2 + 사전 등록 18533d6(00:20): 주 팔 ND-1 \texttt{nd-\allowbreak{}xyz}(RGB만, 격자·탁자 높이 없이 xyz --- 암묵)·ND-2 \texttt{nd-\allowbreak{}est}(탁자 높이·물체 위치·크기를 먼저 추정한 뒤 xyz), 비교 팔 \texttt{pt}(점 0--1000 + 높이 의도, 코드가 머리 깊이·보정·평면 위 군집으로 목표)·L8-xyz(격자), 선택 \texttt{nd-\allowbreak{}pt}. F0 영샷 점, F2 같은 L8 편·816걸음 LoRA $\times$ 팔, DEV + OOD-H(탁자 0.82·0.88). 판정 규칙은 결과 전 고정.
  \item \textbf{결과:} 변경 1(fbf45f7, 00:36, 본 실행 전 스모크 DEV 8편): 깊이 해석기에 로봇 자기 가림, 큰 놓을 곳 라벨 허용 25 mm, 관문 G1 = 글 + 라벨 동일(pt 라벨 결측 4/60, v2 요청 60/60·라벨 60/60이 L8과 동일, 머리 PNG 평균 차 0.46/255). 855ffa5(00:50): 평가 세트에 팔 라벨이 없는 상태도 유지(label\_missing).
  \item \textbf{결과:} 변경 2(64bb9d6, 02:06, 미세조정 팔 답은 아직 없음 --- 학습 중): 이미 본 기준 L8-xyz의 OOD-H($\pm$3 cm)는 오프라인 접근 xy 중앙 2.5 mm·접근 3D 중앙 31.2 mm(z가 탁자 변화만큼 어긋남), 폐루프 4/4(패드 4.5 cm가 $\pm$3 cm z 어긋남을 흡수) --- $\pm$3 cm로는 팔을 가르기 어려워 탁자 0.78·0.92 m($\pm$7 cm) 넓은 OOD-H를 둘째 판정으로 추가(주 판정은 $\pm$3 cm 그대로).
  \item \textbf{진행 중·대기:} \tentative{결과 문서(\texttt{results/\allowbreak{}pt.\allowbreak{}md})는 아직 없다 --- 진행 중. 유료 비교(Astra·gpt-5.2 $\times$ pt/nd-est)는 제안만, 사용자 허락 전 실행 금지.}
\end{itemize}

\paragraph{B8. 최종 상위 = 35B, 8B는 임시 대리}
{\small 09-27 01:2x -- 01:30 KST}
\begin{itemize}
  \item \textbf{사용자} (ul 125, 01:2x KST): \intent{상위를 8B로 하는건 임시인거고 35B짜리는 학습하는게 너무 커야하니까. 최종은 35B가 목적인거야. ㅇㅋ?}
  \item \textbf{그래서(결정):} 정본 §96 보충 2(ce7ba33, 01:30): 최종 상위 = Qwen3.5-35B-A3B, 8B(E-TEACH-L8)는 35B를 바로 학습하기엔 너무 커서 경향만 먼저 보는 임시 대리. 같은 보충의 의도 공유 프로토콜은 C9.
\end{itemize}

\paragraph{B9. 손으로 쓴 프롬프트 정보가 추론 때도 필요한가 $\rightarrow$ E-STRIP8}
{\small 09-27 01:2x -- 01:35 KST}
\begin{itemize}
  \item \textbf{사용자} (ul 126, 01:2x KST): \intent{저 프롬프트는 누가 만들어준건데? 저거는 학습할때는 들어갈만한데 추론시에도 있었다는거?}
  \item \textbf{사실 확인:} \texttt{astra-\allowbreak{}solo@v2} 프롬프트는 우리(Claude 에이전트)가 손으로 썼고, E-TEACH-L8은 학습·추론 모두 이 프롬프트를 썼다: 탁자 높이 0.850, 물체 모양·크기, 그리퍼 기하, 잡기·놓기 요령, 카메라 자세, 작업 상자, 탁자 높이에 그린 격자·낙하선.
  \item \textbf{해 봄:} 사전 등록 5511d5d(01:32, 학습 전, 무료): 팔 S-full(L8 그대로)·S-min(학습·추론 모두 최소 요청: 탁자 높이·격자·낙하선·물체 크기·요령 뺌, 로봇 자기 정보는 남김)·S-drop(학습 때 특권 정보 드롭아웃, 추론은 최소)·S-full$\rightarrow$min(L8 어댑터에 최소 요청)·S-stale(OOD-H에 손으로 쓴 0.850). 같은 L8 상태·816걸음, DEV 305 + OOD-H 오프라인, 판정 NOT\_NEEDED / NEEDED / UNCLEAR은 결과 전 고정. \texttt{strip(v2, \{table, overlay\})} = E-PT \texttt{nd-\allowbreak{}xyz@v1}와 바이트 동일(시험).
  \item \textbf{결과:} 변경 1(8700871, 01:35, 학습·평가 전): G-stale 기준 = 글 동일 + 영상당 다른 화소 $\leq$ 100(재생성 142장 중 48장이 서브 mm TCP 차로 낙하선 1--69 화소가 다름).
  \item \textbf{진행 중·대기:} \tentative{학습·평가 전 --- 결과 없음.}
\end{itemize}

\paragraph{B10. 35B-A3B 바로 할 수 있게 준비만 --- 지연 약 1 s (준비 측정, 과제 성능 아님)}
{\small 09-27 01:0x -- 02:4x KST}
\begin{itemize}
  \item \textbf{사용자} (ul 135, 01:0x KST): \intent{35B-A3B는 일단 바로 할 수 있게 준비만 다 해둬보자}
  \item \textbf{해 봄:} 준비만(본 학습 안 함, 1e4ed20): \texttt{Qwen/\allowbreak{}Qwen3.\allowbreak{}5-\allowbreak{}35B-\allowbreak{}A3B}·FP8판 내려받기, 서빙 지연 실측, 학습 스모크 50걸음, 병합-서빙 왕복, 실행 스크립트 \texttt{tools/\allowbreak{}teach\_\allowbreak{}35b/\allowbreak{}train.\allowbreak{}sh}, 사전 등록 초안 \texttt{docs/\allowbreak{}stage3/\allowbreak{}prereg\_\allowbreak{}teach\_\allowbreak{}35b.\allowbreak{}md}(한계 찾기 지표 3.1 포함).
  \item \textbf{결과:} \texttt{results/\allowbreak{}teach\_\allowbreak{}35b\_\allowbreak{}ready.\allowbreak{}md}: FP8·H200 한 장·한 요청씩·생각 끔 영샷 --- solo p50 1.11 s·p95 1.26 s(유효 59/60), couple@v2 p50 1.50 s; 생각 켬은 6,000토큰까지 잘려 쓸 수 없음. 학습 스모크(fla 커널) 6.07 s/걸음 $\rightarrow$ L8 2 에폭 약 1.4 GPU-h(8B와 같음). 병합 $\rightarrow$ vLLM 왕복 DEV 20요청 모두 유효. 유료 0원.
  \item \textbf{그래서(결정):} 학습·평가 기본 경로는 solo 형식(ul 134·§98, C11) --- 결합 형식 지원은 코드에 남긴다.
\end{itemize}

\paragraph{B11. 명령어 범위 --- 집고 놓기 먼저, 한계 확인 뒤 확장}
{\small 09-27 02:5x KST}
\begin{itemize}
  \item \textbf{사용자} (ul 136, 02:5x KST): \intent{오픈데이터의 경우엔 엄청 다양한 테스크도 있는거지?}
  \item \textbf{사용자} (ul 136, 02:5x KST): \intent{이게 맞지}
  \item \textbf{그래서(결정):} 상위 명령어는 당분간 지금 형식(손끝 목표·edit·그리퍼·stop, 위에서 잡기)으로 두고, 공개 데이터의 다양한 과제(RobotisSW 322종, BEHAVIOR 50과제, RoboTwin 50종, Robo2VLM 3,396과제 등)는 인식·판단 학습에만 쓴다. 상위 단독의 한계가 확인되면(§98) 손목 회전·접근 방향·밀기·당기기로 명령어를 넓힌다. 논문 3.2절에 한 문장으로 반영.
\end{itemize}
